% -----------------------------------------------
% Template for ISMIR LBD Papers
% 2026 version, based on previous ISMIR templates

% Requirements :
% * 2+n page length maximum
% * 10MB maximum file size
% * Copyright note must appear in the bottom left corner of first page
% * Clearer statement about citing own work in anonymized submission
% (see conference website for additional details)
% -----------------------------------------------

\documentclass{article}
\usepackage[T1]{fontenc} % add special characters (e.g., umlaute)
\usepackage[utf8]{inputenc} % set utf-8 as default input encoding
\usepackage[lbd]{ismir}
% !!!PARA VERSAO A SUBMETER USAR LINHA ABAIXO, REMOVER A ACIMA!!!
% \usepackage[lbd,submission]{ismir} 
\usepackage{amsmath,cite,url}
\usepackage{graphicx}
\usepackage{color}
\def\lbd{}	% Flag to use correct LBD settings in the paper, please do not modify this line
\usepackage{listings}
\usepackage{xcolor}
\usepackage{comment}
\definecolor{cmtgray}{gray}{0.42}
\lstdefinestyle{layerit}{%
  language=Python,
  %basicstyle=\ttfamily\scriptsize,      % 7 pt Courier: 55 characters fit one ISMIR column
  %basicstyle=\fontsize{7}{8.4}\selectfont,
  basicstyle=\fontfamily{zi4}\fontsize{7}{8.4}\selectfont,
  %basicstyle=\fontfamily{zi4}\fontsize{6.9}{7.6}\selectfont,
  %basicstyle=\fontfamily{cmtt}\fontsize{7}{8.4}\selectfont,
  commentstyle=\color{cmtgray}\itshape,
  keywordstyle=\bfseries,
  showstringspaces=false,
  columns=fullflexible, keepspaces=true,
  breaklines=false,                     % every line is short enough that it never has to break
  aboveskip=0pt, belowskip=0pt, xleftmargin=0pt, numbers=none}
\lstset{style=layerit}

% Title. Please use IEEE-compliant title case when specifying the title here,
% as it has implications for the copyright notice
% ------
\title{LayerIt: \linebreak Towards a Framework for Time-Aligned, Composable Music Visualisations}

% Note: Please do NOT use \thanks or a \footnote in any of the author markup

% Single address
% To use with only one author or several with the same address
% ---------------
\oneauthor
  {Fernando Azeredo \hspace{1cm} Ant\'onio S\'a Pinto}
  {Faculdade de Engenharia da Universidade do Porto, Portugal \\
   INESC TEC, Porto, Portugal \\
   {\tt up202308655@fe.up.pt}}

% For the author list in the Creative Common license and PDF metadata, please enter author names.
% Please abbreviate the first names of authors and add 'and' between the second to last and last authors.
\def\authorname{F. Azeredo and A. S\'a Pinto}

% Optional: To use hyperref, uncomment the following.
\usepackage[bookmarks=false,pdfauthor={\authorname},pdfsubject={\papersubject},hidelinks]{hyperref}
% Mind the bookmarks=false option; bookmarks are incompatible with ismir.sty.

\sloppy % please retain sloppy command for improved formatting

\begin{document}

%
\maketitle
%
\begin{abstract}
Figures that relate audio, algorithmic output and notation on a single time axis are ordinarily assembled by hand, with the score positioned over the plots by eye: nothing ties its placement to the signal, and nothing survives a change of recording or algorithm. We present LayerIt, an open-source Python library that composes such figures programmatically. Given a score already warped into performance time and a note-level alignment, it places independently generated components on that shared axis and emits a single SVG in which the score's MEI structure and each added component remain identifiable. We demonstrate it on a beat-tracking analysis of Debussy's \emph{Clair de Lune}, in which the notation exposes a metrical-level error that no signal representation reveals.
\end{abstract}
%
\section{Introduction}\label{sec:introduction}

Music can be represented in many ways: as engraved notation, symbolic encodings, or audio recordings of performances. These representations, however, live in different notions of time. Recordings unfold in physical time, measured in seconds or samples; symbolic encodings organize it in abstract units such as beats and measures; and engraved notation renders it as space, in pixels or millimetres. Establishing correspondences between them, a task known as alignment or synchronisation, is central to both music information retrieval (MIR) and digital musicology.

Alignment algorithms, most prominently for the audio-to-score task, must balance accuracy against computational cost and robustness\cite{muller2021SyncToolboxPython}; this alone takes several forms, such as tolerance to performance errors \cite{nakamura2017PerformanceErrorDetection} or to structural mismatches between score and performance \cite{fremerey2010HandlingRepeatsJumps}. Two further axes cut across these choices: where matching takes place, within the audio or symbolic domain, or across the two \cite{mueller2019CrossModalMusicRetrieval}; and the machinery relied on, from classical signal-processing features to learned representations \cite{zeitler2025RobustAccurateAudio}. 

Once computed, an alignment can be put to work in more than one way. One concern is representation: alignments are often stored in ad hoc, single-use formats, and recent modelling relates musical timelines across the graphical, logical, and physical domains so that positions and annotations move between them and the alignment can be reused \cite{hentschel2026TimeAlignModelling}. A separate concern, which this work takes up, is visualization. The prevailing strategy fixes the score's layout and links its elements to instants on a physical timeline through drawn connectors \cite{thomas2012LinkingSheetMusic}, as do interactive tools that track a performance with a moving score cursor \cite{jiang2025PianoPrecisionAdvancing}. Score and physical time thus stay in separate coordinate systems, leaving the user to bridge them mentally.

An alternative dispenses with anchors altogether: the notation itself is laid out in performance time, each element repositioned according to its performed onset, so that the score stays fully readable as notation while doubling as a visual index into physical time \cite{goebl2025LetsScoreWarpAgain}. This opens the possibility that motivates the present work: a shared performance-time axis onto which further material can be layered, from waveforms and feature curves to other information \cite{pinto2020ShiftIfYou}. Of particular interest are annotations such as beat positions, metre, or tempo indications: explicit in musical notation, yet recoverable from audio only by estimation.

What is lacking is a simple, programmatic way to compose these layers reproducibly: such composites are otherwise assembled by hand, the score positioned over the plots by eye, so nothing ties its placement to the signal and nothing survives a change of recording or algorithm. 
We present LayerIt: given a score--performance alignment, it places independently generated components onto this shared axis and merges them into a single, self-contained visualization. Each is handled as an independent layer that can be added, combined, or removed, making such composites straightforward to build and adapt. Implemented as an open-source Python library, LayerIt\footnote{\url{https://github.com/FernandoAze/LayerIt}} composes this material where it is already produced — as activation arrays, spectrograms and onset lists — rather than exporting it to a separate drawing tool.
%
\section{The LayerIt Framework}\label{sec:framework}

A LayerIt figure is a stack of panels that share one performance-time axis.
Four inputs define it: a performance recording, read with librosa \cite{mcfee2015LibrosaAudioMusic}; a score already warped onto that axis by ScoreWarp \cite{goebl2025LetsScoreWarpAgain}; a note-level alignment of onsets to score elements \cite{weigl2020MultimodalMusicInformation}, which LayerIt consumes rather than computes; and any number of algorithmic outputs.
The output is a single SVG document.

All horizontal placement derives from two numbers read from the warped score's time axis: the span it covers in seconds and its length in pixels.
Their ratio is the only scale factor.
Each layer converts its native units (e.g.\ samples, frames) into seconds as it loads.
A shared context then maps seconds to pixels identically for all of them.
Each layer decides its own vertical placement; the horizontal is decided for it.

A representation joins a figure as a class implementing two methods: one that loads and validates its data, one that draws it into a Matplotlib \cite{hunter2007Matplotlib2DGraphics} axes, so that existing plotting knowledge transfers unchanged.
A third, optional method lets a layer emit its own SVG group.
Ten layers are implemented on this interface, standing on four shapes---\texttt{Curve} (a function of time), \texttt{Events} (instants), \texttt{Intervals}, and \texttt{Field} (a dense time--frequency array)---and none of them touches the aligning machinery.

The score is copied into the composite rather than redrawn from its Verovio rendering \cite{pugin2014VerovioLibraryEngraving}, so its Music Encoding Initiative (MEI) identifiers and hierarchy are preserved.
Each added component is emitted as its own named group, with identifiers rewritten per panel so that repeated layers stay distinguishable; the vector shapes remain native SVG primitives, selectable and restylable, while a field embeds a raster image inside its group.

\section{Demonstration}\label{sec:demonstration}

We demonstrate LayerIt in the figure a beat-tracking analysis calls for: a
recording, a score, and a tracker's frame-level output. The recording and the tracker's output already lie in physical time; the score does not, and is warped into it by ScoreWarp from an onset alignment, so its horizontal placement is computed rather than set by hand. The material is the first six bars of Debussy's \emph{Clair de Lune}, performed by Maria João Pires; the onsets were aligned to the MEI score in Piano Precision \cite{jiang2025PianoPrecisionAdvancing}, and the beat and downbeat activations and estimates were obtained from Beat This!~\cite{foscarin2024BeatThisAccurate}.

Listing~\ref{lst:script} shows the complete script that produces \figref{fig:composite}: five statements in all. The horizontal axis is established once from the warped score; panel order and relative panel height are the only placements specified by the script.
\begin{lstlisting}
fig = Visualizer(audio, score, maps, beats)
fig.add_panel(Onset(), BeatLogits(), DownbeatLogits(), height=.5)
fig.add_panel(Waveform(), BeatsLayer(), DownbeatsLayer(), height=.5)
fig.add_panel(MelSpec(freq_window=(100, 1500)))
fig.compose("LDB_FIG.svg")
\end{lstlisting}


\begin{comment}
\begin{lstlisting}[
  float=t,
  caption={Abridged script generating Figure~\ref{fig:composite}.},
  label={lst:script}
]
fig = Visualizer(audio=audio, score=score, maps=maps, beats=beats)
fig.add_panel(Onset(), BeatLogits(), DownbeatLogits(), height_scale=.5)
fig.add_panel(Waveform(), BeatsLayer(), DownbeatsLayer(), height_scale=.5)
fig.add_panel(MelSpec(freq_window=(100, 1500)))
fig.compose("LDB_FIG.svg")
\end{lstlisting}
\end{comment}

\begin{comment}
\begin{figure}[t]
\begin{lstlisting}
fig = Visualizer(audio=audio, score=score, maps=maps, beats=beats)
fig.add_panel(Onset(...), BeatLogits(...), DownbeatLogits(...),
              height_scale=0.5)
fig.add_panel(Waveform(...), BeatsLayer(...), DownbeatsLayer(...),
              height_scale=0.5)
fig.add_panel(MelSpec(freq_window=(100, 1500)))
fig.compose("LDB_FIG.svg")
\end{lstlisting}
\caption{The script that generates Figure~\ref{fig:composite}: the inputs declared once, one \texttt{add\_panel} call per panel, and a single \texttt{compose} call.}
\label{fig:script}
\end{figure}
\end{comment}

\begin{figure}
\includegraphics[width=\columnwidth]{output.png}
\caption{LayerIt composition, top to bottom: warped score (1); beat/downbeat logits (2) and estimates (3), with beats in dotted orange, downbeats in solid blue, and score onsets as dashed vertical lines; mel spectrogram (4).}
\label{fig:composite}
\end{figure}

Most downbeats are correctly placed, but the many false positives---several of them quaver onsets taken as beats---show that the metrical structure is misread. Two difficulties are visible at once: compound triple metre is scarce in the canonical beat-tracking datasets, and the playing is markedly expressive, as the \emph{Andante très expressif} marking announces and the warped score confirms, its spacing stretching and compressing the dotted crotchet throughout. Read against the notation, the estimates become an account of the difficulties behind them rather than a list of errors.

%Substituting a different beat tracker changes one path in the script; a different performance additionally requires its own alignment and warped score, which LayerIt does not compute.

\section{Conclusion}\label{sec:conclusion}

% LayerIt provides a programmatic workflow for generating score-aligned MIR figures from independently specified components. Its structured SVG output preserves the score's semantic information and the identity of added components, rather than flattening the figure into a single raster image. The current implementation is open source and provides a stable core for additional layer types. Planned extensions include tempograms and spectrogram derivatives, time-alignment backends beyond \textit{ScoreWarp}, and annotation export. Integrating an upstream synchronisation toolkit such as \textit{Sync Toolbox}, and supporting interchange data such as match files, would broaden the set of alignment inputs available to the framework. The structured SVG representation may also support interactive, playback-synchronised viewers. By making the figure specification re-runnable, LayerIt supports transparent and reusable visualisation workflows.

LayerIt composes score-aligned MIR figures from independently specified components, and its structured SVG preserves the score's semantic information and the identity of each added component rather than flattening the figure into a raster image. We demonstrated that this is more than a convenience: a diagnosis that depends on notated structure---metre here, but equally repeats, voicing or articulation---is unavailable in a signal-only view, however many signal representations are stacked.

The framework depends on an externally computed alignment and inherits its errors, and dense time--frequency representations remain the one component that is rasterised.
%As LayerIt operates on a precomputed score–performance alignment, any errors in that alignment propagate directly to the resulting visualisation. Such errors remain a known issue in score–audio alignment and transcription, particularly for structurally complex performances, difficult recording conditions, and polyphonic or orchestral material.

Two directions follow. The first concerns the alignment input. LayerIt reads the onset-level format that ScoreWarp consumes; reading alignments produced by the Sync Toolbox \cite{muller2021SyncToolboxPython} would extend it to automatic synchronisation, and adopting a structured representation such as Time to Align! \cite{hentschel2026TimeAlignModelling} would preserve the provenance and domain information that an onset list does not carry. The second is longer term. The same composite is a natural interface for annotation rather than inspection: beat annotation is normally carried out against audio alone, with the metrical level inferred from the signal, whereas a score-aligned view makes it explicit. Since the output is a structured SVG in which every component remains addressable, an interactive, playback-synchronised annotation interface is an extension of the present design rather than a departure from it.

%\newpage

%\section{AI Usage Statement}

%Large language models (LLMs) were used in this work as a writing and code development aid. Specifically: (1)~LLMs were used to help structure and write this LBD paper, including organization, phrasing, and clarity improvements; (2)~LLMs were used to assist in the development of the LayerIt codebase. The core research methodology, conceptual contributions, framework architecture, and scientific validation remain the work of the authors. All code, design decisions, and experimental results have been authored and verified by the authors themselves.

% For bibtex users:
\bibliography{ISMIRtemplate_ASP}

\end{document}
